\chapter{Stand der Technik}

\section{Tor}
//TODO: Tor Client auch als CLI Tool verfügbar

Das Tor Projekt stellt den Tor Browser zur Verfügung.
Dieser ist ein Fork von Firefox und wurde für die Nutzung mit Tor optimiert.
Der Tor Browser greift auf den eigentlichen Tor Client zu und leitet den gesamten Internetverkehr des Browsers über einen SOCKS5 Proxy in das Tor Netzwerk weiter.
Der Tor Client leitet den Datenverkehr über langlebige kryptographische Tunnels, Circuits genannt, weiter.\cite[S. 2]{jansen2023data}

Jeder Circuit enthält 3 Tor Knoten. Den Eingangsknoten, Mittelknoten und den Ausgangsknoten.
Der Eingangsknoten kennt den Client und den Mittelknoten nicht aber den Ausgangsnoten und das Ziel.
Der Mittelknoten kennt den Eingangs und Ausgangsknoten nicht aber den Client oder das Ziel.
Und der Ausgangsknoten letztendlich kennt lediglich den Mittelknoten und das Ziel, nicht aber den Client.
Somit wird die Anonymisierung des Internetverkehrs erzeugt.

Ein Tor Circuit lebt für etwa 10 Minuten und ist relativ Teuer in der Erstellung (Diffie Hellmann Key Exchange.)

Ein Tor Stream ist eine TCP Verbindung welche über einen Tor Circuit geleited wird.

Falls eine Anwendung mehrere Streams gleichzeitig versucht Tor diese alle über einen Circuit zu leiten.
So kann unter anderem die Chance minimiert werden über einen bösartigen Knoten geleitet zu werden.
Im Beispiel des Tor Browsers werden jeweils alle Streams einer Top-Level Domain über einen Circuit geleitet. \cite[S. 2]{jansen2023data}

Der Tor Datenverkehr wird dann in 514 Byte grosse Pakete verpackt, welche auch Tor Zellen genannt werden.
Dadurch dass die Pakete immer genau gleich gross sind soll die Überwachung vom Internetverkehr weiter erschwert werden.
Die Anzahl der Zellen und deren Richtung kann von einem Angreifer jedoch immernoch ausgelesen werden.
In einer Website Fingerprinting Attacke wird genau dies ausgenützt. \cite[S. 2]{jansen2023data}

\section{Onion Routing}
//TODO: Onion Routing beschreiben



\section{Website Fingerprinting}

In einem \gls{wf} Angriff kann der Angreifer die Metadaten des Internetverkehrs vom Opfer auslesen.
Dabei wird der Datenverkehr zwischen dem Client und dem Eingangsknoten überwacht.
Der Angriff kann dabei sowohl vom Betreiber vom Eingangsknosten\cite[S. 2]{jansen2023data} oder auch von einem passiven Beobachter durchgeführt werden.\cite[S. 3]{wang2016realistically}

Der Angreifer hat eine vordefinierte Liste an Websiten die er kontrollieren möchte.
Er klassifiziert nun den überwachten Datenverkehr ob eine der definierten Websiten aufgerufen wird oder nicht.
Der \gls{wf} Angriff gilt als Erfolgreich wenn ein Aufruf korrekt erkannt wurde, aber auch wenn ein Nicht-Aufruf korrekt erkannt worden ist.

Der Angreifer eine Menge an überwachten Webseiten (\textit{monitored pages}), deren Zugriffe er erkennen möchte.
Alle übrigen Webseiten gelten als nicht-überwacht (\textit{non-monitored}).
Der Klassifikator entscheidet für jede beobachtete Zellsequenz, ob sie einer überwachten Seite zuzuordnen ist (positiv) oder nicht (negativ).

Der \gls{ml} Klassfikator wird dabei für zwei Aufgaben trainiert.
Bei der binären Klassifizierung soll der Klassifikator mitteilen ob der Aufruf zu einer überwachten Seite führt oder nicht.
In einem zweiten Schritt wird eine Multiclass Klassifizierung vorgenommen.
Es soll erkannt werden um welche Seite es sich handelt falls der erste Schritt in einem positiven Ergebnis mündete. \cite[S. 3]{jansen2023data}

Dabei ergeben sich vier mögliche Ausgänge:
\begin{itemize}
    \item \textbf{True Positive (TP):} Eine überwachte Seite wird korrekt als solche erkannt.
    \item \textbf{True Negative (TN):} Eine nicht-überwachte Seite wird korrekt als
          nicht-überwacht eingestuft.
    \item \textbf{False Positive (FP):} Eine nicht-überwachte Seite wird fälschlicherweise
          als überwacht klassifiziert.
    \item \textbf{False Negative (FN):} Eine überwachte Seite wird nicht erkannt.
\end{itemize} \cite[S. 3]{jansen2023data}.

Die \textit{True Positive Rate} (TPR) gibt den Anteil korrekt erkannter Zugriffe auf überwachte Seiten an, während die \textit{False Positive Rate} (FPR) den Anteil fälschlicherweise als überwacht eingestufter Zugriffe misst:
\begin{equation}
    \text{TPR} = \frac{\text{TP}}{\text{TP} + \text{FN}}, \qquad
    \text{FPR} = \frac{\text{FP}}{\text{FP} + \text{TN}}
\end{equation} \cite[S. 3]{jansen2023data}
22
Eine niedrige FPR ist besonders entscheidend, da die \textit{Base Rate}, also die Wahrscheinlichkeit, dass ein Nutzer tatsächlich eine überwachte Seite besucht, in der Praxis sehr gering ist.
Bei hoher FPR würden die meisten positiven Klassifikationen Fehlalarme sein, selbst wenn die TPR hoch ist (\textit{Base Rate Fallacy}) \cite[S. 3]{wang2016realistically}.




Darüber hinaus wurde die Simulationstreue von Shadow evaluiert.
Die Netzwerk-Performance der simulierten Tor-Umgebung stimmte in den Median-Werten weitgehend mit dem Live-Netzwerk überein, wobei Abweichungen vor allem in den Long-Tail-Bereichen auftraten. 
In einem Cross-Dataset-Experiment erreichte ein ausschliesslich auf Shadow-Daten trainierter Klassifikator eine Accuracy von über 85\,\% auf Live-Tor-Traces, was die Eignung von Shadow für kontrollierte Website-Fingerprinting-Evaluationen bestätigt. \cite[S. 6-8]{jansen2023data}

In den Abschnitten zur Sensitivität und Robustheit untersuchten Jansen und Wails, wie empfindlich WF-Klassifikatoren auf veränderte Netzwerkbedingungen reagieren. 
Die Zellsequenzen wurden dabei an den Entry Relays der simulierten Tor-Netzwerke aufgezeichnet, was dem realistischen Bedrohungsmodell eines lokalen Angreifers entspricht. \cite[S.~8--10]{jansen2023data}

Es zeigte sich, dass die Closed-World-Accuracy weitgehend stabil blieb, während die False-Positive-Rate im Open-World-Szenario um bis zu 700\,\% anstieg, wenn Trainings- und Testdaten aus unterschiedlich belasteten Netzwerken stammten. 
Insbesondere zeitbasierte Klassifikatoren wie Tik-Tok und $k$-Fingerprinting erwiesen sich als besonders anfällig gegenüber variierenden Netzwerkbedingungen. 
Durch ein robustes Trainingsverfahren, bei dem Trainingsdaten aus mehreren Netzwerkkonfigurationen gemischt wurden, konnte die FPR um bis zu 83\,\% gesenkt werden. \cite[S.~8--10]{jansen2023data}

Als Klassifikator wird Deep Fingerprinting (DF) nach Sirinam et al. \cite{sirinam2018deep} verwendet.
DF basiert auf einem Convolutional Neural Network, das ausschliesslich Zellrichtungen als Eingabe verwendet und somit zeitunabhängig arbeitet.
Dies ermöglicht eine isolierte Betrachtung des Einflusses von Circuit Padding auf die richtungsbasierten Verkehrsmuster, ohne dass simulationsbedingte Timing-Variationen die Ergebnisse verfälschen.\cite{jansen2011shadow}
//TODO: Erwähnen und beschreiben wie DF genutzt und validiert wurde in dem paper von Jansen und wails 




Es gibt öffentlich verfügbare Dateien im ZIM Format welche die vollständige Wikipedia Sammlung beeinhalten.

Der Tor Browser selbst kann im Shadow Simulator nicht ausgeführt werden.
Deshalb greifen wir auf wget2 zurück

\footnote{\url{https://github.com/rockdaboot/wget2}}





\subsection{Website-Fingerprinting-Klassifikatoren}

In der WF-Literatur haben sich verschiedene Klassifikatoren etabliert, die sich 
entlang zweier Dimensionen unterscheiden lassen: klassische 
Machine-Learning-Verfahren gegenüber neuronalen Netzen sowie zeitbasierte 
gegenüber zeitunabhängigen Ansätzen. Im Folgenden werden vier repräsentative 
Klassifikatoren vorgestellt, die von Jansen und Wails in ihrer Evaluierung 
verwendet wurden \cite[S.~9\,f.]{jansen2023data}.

\paragraph{CUMUL}
CUMUL von Panchenko et al. \cite{panchenko2016website} ist ein klassischer, 
zeitunabhängiger Ansatz. Der Klassifikator verwendet eine Support Vector Machine 
(SVM) mit Radial-Basis-Funktions-Kernel. Als Eingabe dient eine kumulativ 
aufsummierte Zellrichtungssequenz $\langle c_i \rangle$ mit 
$c_i = \sum_{j=0}^{i} d_j$, wobei $d_j \in \{-1, 1\}$ die Richtung der $j$-ten 
Zelle angibt. Die kumulative Sequenz wird anschliessend durch lineare Skalierung 
auf genau 100 Elemente normiert.

\paragraph{$k$-Fingerprinting}
$k$-Fingerprinting ($k$-FP) von Hayes und Danezis \cite{hayes2016kfingerprinting} 
kombiniert einen Random-Forest-Klassifikator mit einem 
$k$-Nearest-Neighbors-Verfahren. Im Multiclass-Setting wird ausschliesslich der 
Random Forest verwendet, der auf Basis von Summary-Statistiken einer Zellsequenz 
klassifiziert -- darunter Minimum, Maximum, Mittelwert und Standardabweichung der 
Zell-Interarrival-Zeiten. Da $k$-FP Timing-Informationen als Features nutzt, 
zählt er zu den zeitbasierten Klassifikatoren.

\paragraph{Deep Fingerprinting}
Deep Fingerprinting (DF) von Sirinam et al. \cite{sirinam2018deep} basiert auf 
einem tiefen Convolutional Neural Network (CNN), das aus mehreren gestapelten 
Blöcken von Convolutional-, Batch-Normalization-, Aktivierungs-, Pooling- und 
Dropout-Schichten besteht. Als Eingabe dient direkt der Zellrichtungsvektor 
$\langle d_i \rangle_{i=1}^{5000}$ mit $d_i \in \{-1, 1\}$; kürzere Sequenzen 
werden mit Nullen aufgefüllt. DF arbeitet zeitunabhängig und hat in früheren 
Arbeiten sowohl CUMUL als auch $k$-FP in der Klassifikationsgenauigkeit 
übertroffen \cite{sirinam2018deep}.

\paragraph{Tik-Tok}
Tik-Tok (TT) von Rahman et al. \cite{rahman2020tiktok} verwendet dieselbe 
CNN-Architektur wie DF, erweitert jedoch die Eingabe um Timing-Informationen. 
Anstelle reiner Zellrichtungen nutzt TT das Produkt aus Richtung und Zeitstempel 
$\langle t_i \cdot d_i \rangle_{i=1}^{5000}$. Damit verwendet TT strikt mehr 
Information als DF. Allerdings haben Jansen und Wails gezeigt, dass zeitbasierte 
Klassifikatoren wie TT empfindlicher auf variierende Netzwerkbedingungen reagieren 
\cite[S.~10]{jansen2023data}.

\paragraph{Wahl des Klassifikators}
In der vorliegenden Arbeit wird Deep Fingerprinting als Klassifikator verwendet. 
Diese Wahl ist durch mehrere Faktoren motiviert: Erstens arbeitet DF 
zeitunabhängig, wodurch simulationsbedingte Timing-Variationen die Ergebnisse 
nicht verfälschen. Zweitens beeinflusst Circuit Padding primär die 
Zellrichtungssequenzen durch das Einfügen zusätzlicher Dummy-Zellen, was direkt 
auf die von DF genutzten Features wirkt. Drittens wurde DF bereits von Jansen und 
Wails erfolgreich in Shadow validiert und erreichte im Cross-Dataset-Experiment 
eine Accuracy von 82\,\% auf Live-Tor-Daten \cite[S.~12]{jansen2023data}. Dies 
gewährleistet die Vergleichbarkeit der vorliegenden Ergebnisse mit der bestehenden 
Literatur.




\subsection{Closed-World- und Open-World-Szenarien}

In der Website-Fingerprinting-Literatur werden zwei grundlegende 
Evaluationsszenarien unterschieden: das Closed-World- und das 
Open-World-Szenario \cite{wang2016realistic, jansen2023data}.

Im \textbf{Closed-World-Szenario} wird angenommen, dass der Nutzer 
ausschliesslich Webseiten aus einer dem Angreifer bekannten, 
endlichen Menge $W_\alpha$ besucht. Der Klassifikator muss 
lediglich entscheiden, welche der bekannten Seiten besucht wurde. 
Da alle Klassen gleich häufig im Testdatensatz vertreten sind, 
wird die Klassifikationsgenauigkeit (Accuracy) als Performanzmass 
verwendet. Dieses Szenario eignet sich gut für kontrollierte 
Experimente, ist jedoch unrealistisch, da es die Existenz 
unbekannter Webseiten vollständig ausschliesst 
\cite[S.~21]{wang2016realistic}.

Im \textbf{Open-World-Szenario} hingegen definiert der Angreifer 
eine Menge sensitiver Seiten $W_\alpha$, die er erkennen möchte, 
während der Nutzer beliebige weitere Seiten besuchen kann -- 
sowohl dem Angreifer bekannte, nicht-sensitive Seiten $W_\beta$ als 
auch völlig unbekannte Seiten $W_\emptyset$. Der Klassifikator 
muss somit zunächst eine binäre Entscheidung treffen: Gehört eine 
beobachtete Zellsequenz zu einer sensitiven Seite oder nicht? Da 
die Klassen in diesem Szenario stark unbalanciert sind -- die 
Anzahl nicht-sensitiver Seiten übersteigt die der sensitiven um 
ein Vielfaches -- ist Accuracy kein geeignetes Mass. Stattdessen 
werden die True Positive Rate (TPR) und die False Positive Rate 
(FPR) herangezogen \cite[S.~3]{jansen2023data}. Eine niedrige FPR 
ist dabei besonders wichtig, da bei der geringen Basisrate 
sensitiver Seitenbesuche selbst eine scheinbar kleine FPR zu einer 
grossen Anzahl falsch-positiver Alarme führen kann 
\cite[S.~23]{wang2016realistic}.

Wang und Goldberg betonen, dass frühere WF-Arbeiten häufig unter 
vereinfachenden Closed-World-Annahmen evaluiert wurden, was die 
tatsächliche Bedrohung durch WF-Angriffe nur unzureichend 
widerspiegelt \cite{wang2016realistic}. Jansen und Wails haben 
jedoch gezeigt, dass die Closed-World-Accuracy selbst bei 
variierenden Netzwerkbedingungen stabil bleibt (Schwankungen von 
nur 1--3\,\%), während TPR und FPR im Open-World-Szenario 
deutlich empfindlicher auf veränderte Bedingungen reagieren 
\cite[S.~10\,f.]{jansen2023data}.

Die vorliegende Arbeit verwendet ein Closed-World-Szenario mit 
einer einzelnen Netzwerkkonfiguration. Diese Wahl ist dadurch 
motiviert, dass der Fokus der Untersuchung auf dem Einfluss von 
Circuit Padding auf die Klassifikationsgenauigkeit von Deep 
Fingerprinting liegt und nicht auf der Robustheit gegenüber 
variierenden Netzwerkbedingungen. Die Erweiterung auf ein 
Open-World-Szenario bleibt zukünftiger Forschung vorbehalten.



\paragraph{Anpassung von wget2}
Da Tor seinen Dienst über einen SOCKS5-Proxy bereitstellt, 
Webcrawler wie \texttt{wget2} jedoch standardmässig nur 
HTTP-Proxies unterstützen, erweiterten Jansen und Wails 
\texttt{wget2} um einen SOCKS-Client. Erst diese Modifikation 
ermöglichte es, \texttt{wget2}-Anfragen über den lokalen 
Tor-Prozess zu leiten und somit Zellsequenzen am Entry Relay 
aufzuzeichnen \cite[S.~5]{jansen2023data}. Darüber hinaus wurde 
\texttt{wget2} so konfiguriert, dass es alle eingebetteten 
Ressourcen einer Seite parallel herunterlädt 
(\texttt{-{}-page-requisites}, \texttt{-{}-max-threads=30}) und 
denselben User-Agent wie der Tor Browser verwendet. Durch lokale 
Tests wurde verifiziert, dass beide Tools -- \texttt{wget2} und 
Tor-Browser-Selenium (TBS) -- einen identischen Satz an 
Ressourcen vom Server abrufen. Die gepatchte \texttt{wget2}-Version 
ist im Artefakt-Repository der Autoren öffentlich verfügbar.